% Shared concise source. Full proof provenance is in the gauss-book source map.
\section{An arithmetic Gauss tail with the original finite cutoff}
\label{gtl:core}
The coherent kernel of Section~\ref{cbk:core} contains a difficult
root-weighted mean. A large part of its \emph{short-divisor expansion}
can now be bounded by retaining a cubic interaction before taking a
triangle. This gives a signed cutoff comparison, not a smaller norm
for the original coefficients. A direct exact-discrete-frequency
version is included as well. The remaining small-short-divisor
moment and the separate discrete-frequency error are stated below.
The whole signed power $13/12$ and global unsigned power $9/8$ do not
change.

\subsection{Original objects and a fixed common profile}
\label{gtl:objects}
Keep $K=\mathbb Q(\omega)$, $\mathcal O=\mathbb Z[\omega]$,
$\delta=1+2\omega$, and the fixed exclusions $S$ of
Section~\ref{cbk:objects}. A good ideal is prime to $S$.
All columns are the complete good squarefree ideal annulus
$D\le N\mathfrak m\le2D$, $D>1$. Equal norms do not identify ideals.
For a real cutoff $1\le T\le\sqrt D$ put
\begin{equation}\label{gtl:coefficients}
 \begin{aligned}
 A_T(\mathfrak m)
 &=\sum_{\substack{\mathfrak a\mathfrak b\mid\mathfrak m\\
             (\mathfrak a,\mathfrak b)=1\\
             N\mathfrak a,N\mathfrak b\le T}}
                    \mu(\mathfrak a)\mu(\mathfrak b),\\
 c_f(\mathfrak m)
 &=f(\mathfrak m)\nu(\mathfrak m)R_{\mathfrak m}(u_0)
                                             W(N\mathfrak m/D).
 \end{aligned}
\end{equation}
Here $\nu$ is the original fixed unitary finite-order character;
$u_0$ is the fixed literal unit/$S$-supported factor with no good
support. The actual coefficient is $c_{A_U}$, $1\le U\le\sqrt D$.
The function $W$ is one fixed, possibly complex, smooth common
profile supported on $[1,2]$, zero at the endpoints. Define
\begin{equation}\label{gtl:profile}
 \begin{aligned}
 B_W&=\|W\|_\infty,\qquad
 \widehat W(t)=\int_0^\infty W(x)x^{-it}\,\frac{dx}{x},\\
 \mathcal M(W)&=\frac1{2\pi}\int_{\mathbb R}|\widehat W(t)|\,dt.
 \end{aligned}
\end{equation}
This explicit seminorm is finite and controls $B_W$. It is a
seminorm of the column profile, not of the sharp nodal row function.

Only the original unweighted \emph{pure} two-pivot block is used:
$A_2=1$, $\mathfrak v_2=\mathcal O$, and the other four factors
are units. The first row is good and squarefree, with the original
literal/ray conditions and capacity. Its integer endpoints are
\begin{equation}\label{gtl:capacity}
 a=\lceil A_1\rceil,\qquad
 b=\min\{\lceil2A_1\rceil-1,\lfloor H/Nu_0\rfloor\},
 \qquad K=\sqrt b .
\end{equation}
An empty interval contributes zero throughout. Otherwise $b\ge a\ge1$.
Use the exact $G$ of~\eqref{cbk:nodal}, its transform
$\dot G$ in~\eqref{cbk:bessel}, and $C_0=3/(4\pi^2)$.
In particular $0\le G\le1$, $\operatorname{supp}G\subset[1/2,5/2]$,
and $\int G\le2$. The original \emph{pre-$S$} cutoff $K$ is never
recomputed after another deletion.

Let $\mathcal I$ be the original ordered coprime column pairs.
One may also impose the \emph{original} unequal-norm projection;
the same choice is then used in every expression below. For each
source pair $A,B$ and original ray character $\rho$ use the primitive
$\eta$, $\mathfrak f_S$ and $F$ of~\eqref{cbk:eta}. Thus
$F=N\mathfrak f_S N\mathfrak mN\mathfrak n$, and $\eta$ has total
infinity type zero. Put
\begin{equation}\label{gtl:deletions}
 \begin{aligned}
 T_q&=S\cup\operatorname{supp}(\mathfrak m\mathfrak n),\\
 Q_s^{T_q}(K;\eta^2)
   &=\sum_{\substack{\mathfrak k\ {\rm good,SF}\\
                   N\mathfrak k\le K,\;(\mathfrak k,T_q)=1}}
         \frac{\mu(\mathfrak k)\eta(\mathfrak k)^2}
              {(N\mathfrak k)^s},\\
 P_S(\eta)&=\sum_{\mathfrak e\mid P_S}
                   \frac{\mu(\mathfrak e)\eta(\mathfrak e)}
                        {N\mathfrak e},
 \qquad P_S=\prod_{p\in S}p .
 \end{aligned}
\end{equation}
The symbol $(\mathfrak k,T_q)=1$ means avoidance of every prime in
that set. Squaring $\eta$ never restores a deleted prime. Values at
$\mathfrak e$ and at any later $S$-supported frequency are separate
primitive values, including zero.

\subsection{The actual retained mean and its finite leading model}
\label{gtl:mean}
In the exact primitive kernel~\eqref{cbk:exact-kernel}, write each
dual ideal uniquely as $\mathfrak y=\mathfrak h\mathfrak r\mathfrak s^6$:
$\mathfrak h$ is $S$-supported, $\mathfrak r$ is good and sixth-free,
and $\mathfrak s$ is good. Overlap of $\mathfrak r,\mathfrak s$
is allowed. For $X\ge64$ set
\begin{equation}\label{gtl:top-window}
 V=X^{1/6}/2,\qquad
 b_{\mathfrak w}=(X/N\mathfrak w)^{1/6},\qquad
 t_{\mathfrak w}=\bigl((N\mathfrak w)^{-1/6}-1/2\bigr)_+ .
\end{equation}
The top band is $V<N\mathfrak s\le2V$, $N\mathfrak y\le X$.
Replacing its $\mathfrak s$ count by its principal fixed-ray
density defines the retained mean $M_{\rm top}$. This operation
keeps, rather than deletes, its coherent main term.
The selected branches are $\rho^6=1$. At this top window
$N\mathfrak hN\mathfrak r<64<7^6$, so
$\mathfrak w=\mathfrak h\mathfrak r$ runs over \emph{all} integral
ideals of norm below $64$, including $S$-supported ideals.

Set
\begin{equation}\label{gtl:density}
 \kappa_S=\frac{\pi}{3\sqrt3}\prod_{p\in S}(1-(Np)^{-1}),
 \qquad
 g(\mathfrak m)=\prod_{p\mid\mathfrak m}(1-(Np)^{-1}).
\end{equation}
On original coprime support the density is
$\kappa_Sg(\mathfrak m)g(\mathfrak n)$.
For $\theta=(A,B,\rho,\mathfrak m,\mathfrak n,\mathfrak e,\mathfrak w)$
in the preceding sets, with $\rho^6=1$, define
\begin{equation}\label{gtl:outer-coefficient}
 \begin{aligned}
 d_\theta[f_1,f_2]
 &=\frac{\kappa_S A_1}{D}h_{AB\rho}(\mathcal O)
     c_{f_1}(\mathfrak m)\overline{c_{f_2}(\mathfrak n)}
     \frac{g(\mathfrak m)g(\mathfrak n)\epsilon(\eta)}{\sqrt F}\\
 &\quad\cdot
      \frac{\mu(\mathfrak e)\eta(\mathfrak e)}{N\mathfrak e}
                       \overline{\eta(\mathfrak w)},\qquad
 \lambda_\theta=\frac{A_1N\mathfrak w}{N\mathfrak e F}.
 \end{aligned}
\end{equation}
The exact definition of the normalized retained mean is
\begin{equation}\label{gtl:actual-mean}
 \begin{aligned}
 \mathcal M_{\rm top}
 :=\frac{M_{\rm top}}D
 &=\sum_\theta d_\theta[A_U,A_U]
     \sum_{\substack{\mathfrak k\ {\rm good,SF}\\N\mathfrak k\le K,\;(\mathfrak k,T_q)=1}}
       \frac{\mu(\mathfrak k)\eta(\mathfrak k)^2}
            {(N\mathfrak k)^2}\\
 &\quad\cdot\int_V^{b_{\mathfrak w}}
       \dot G\left(\frac{\lambda_\theta t^6}
                         {(N\mathfrak k)^2}\right)\,dt .
 \end{aligned}
\end{equation}
All sums here are finite. In particular there is no
$\mathfrak k\mid\mathfrak w$ or $(\mathfrak k,\mathfrak w)=1$
condition. The finite original source/ray coefficients
$h_{AB\rho}(\mathcal O)$ have not been centered.
An argument at $\mathfrak e$ is not divided by a value at
$\mathfrak w$.

Put $g_0=C_0^{-1/2}\int G$ and define
\begin{equation}\label{gtl:leading-model}
 \mathcal L_0[f_1,f_2]
  =g_0X^{1/6}
      \sum_\theta d_\theta[f_1,f_2]\,t_{\mathfrak w}
                             Q_2^{T_q}(K;\eta^2).
\end{equation}
Replacing both column functions here is an explicitly defined
auxiliary model, not a change to the original primal data.

\begin{proposition}[Matched leading model with its original cutoff]
\label{gtl:model}
For the actual coefficients and every finite $X\ge64$,
\begin{equation}\label{gtl:model-error}
 |\mathcal M_{\rm top}-\mathcal L_0[A_U,A_U]|
 \ll B_W^2(DX)^\epsilon\,\frac{A_1^2X^{7/6}}{D^2}.
\end{equation}
The constant uses no derivative bound for the sharp $G$.
Both interval endpoints and the original finite $K$ occur in
$\mathcal L_0$. This estimates the model error, not the leading
signed model.
\end{proposition}

\subsection{The root is an incomplete cubic Gauss weight}
\label{gtl:root}
For each good ideal use its primary generator
$\alpha_{\mathfrak a}\equiv1\pmod3$ and write
$\chi_{j,\mathfrak a}$ for its positive residue character of order
$j=3,6$. For good squarefree $\mathfrak m$ put
\begin{equation}\label{gtl:cubic-gauss}
 \begin{aligned}
 G_{\mathfrak m}
 &=\frac1{\sqrt{N\mathfrak m}}
   \sum_{x\bmod\alpha_{\mathfrak m}}
      \chi_{3,\mathfrak m}(x)
      \exp\!\left(2\pi i\operatorname{Tr}
                   \frac{x}{\delta\alpha_{\mathfrak m}}\right),\\
 g_{\mathfrak m}(z)
 &=\overline{\chi_{3,\mathfrak m}(z)}G_{\mathfrak m}.
 \end{aligned}
\end{equation}
The second expression is the normalized finite Gauss sum at
frequency $z$, including zero when $(\mathfrak m,z)>1$.
For $\mathfrak q=\mathfrak m\mathfrak n$ put
$\theta_q=\chi_{6,\mathfrak m}\overline{\chi_{6,\mathfrak n}}$.
Let $\theta_S$ be the \emph{actual finite $S$ part} of the
total type-zero $\eta$, at its own primitive conductor
$\mathfrak f_S=(\gamma_S)$. Then
\begin{equation}\label{gtl:S-factor}
 \mathcal S_{\mathfrak m,\mathfrak n,\rho}
  =\theta_S(\alpha_{\mathfrak m}\alpha_{\mathfrak n})
       \theta_q(\gamma_S)
        \frac{\tau_{\gamma_S}(\theta_S)}{\sqrt{N\mathfrak f_S}} .
\end{equation}
The unit-conductor convention is one. This scalar has modulus one;
it is not obtained by discarding the infinity type of a constituent.
The accepted Jacobi character in~\eqref{sw:character} gives
$\vartheta=\Psi/\sqrt N$, with conductor $(2\delta)$ of norm $12$
and nonzero infinity type. The exact root formula is
\begin{equation}\label{gtl:root-cubic}
 \begin{aligned}
 \epsilon(\eta)
 &=\mathcal S_{\mathfrak m,\mathfrak n,\rho}
       \Gamma(\alpha_{\mathfrak m}\alpha_{\mathfrak n})
       \mu(\mathfrak m)\mu(\mathfrak n)
       \vartheta(\mathfrak m)\overline{\vartheta(\mathfrak n)}
       \overline{G_{\mathfrak m}}G_{\mathfrak n},\\
 \Gamma(a+b\omega)
 &=\tfrac12\{1+i^{-b}+i^a+i^{b-a}\}.
 \end{aligned}
\end{equation}
No extra quadratic sign is inserted. In the original kernel,
$\eta(\mathfrak k^2)$ and the second-row phase, when present,
produce the raw cubic frequency
$\overline{g_{\mathfrak m}(\alpha_{\mathfrak k}\alpha_{\mathfrak v_2})}
 g_{\mathfrak n}(\alpha_{\mathfrak k}\alpha_{\mathfrak v_2})$;
here $\mathfrak v_2=\mathcal O$. This incomplete weighted
frequency set is not a complete residue average.

\subsection{A signed short-divisor tail from a physical cubic matrix}
\label{gtl:tail}
For squarefree $\mathfrak m$ there is an exact sign-preserving
assignment identity:
\begin{equation}\label{gtl:ternary}
 \mu(\mathfrak m)A_T(\mathfrak m)
   =\sum_{\substack{\mathfrak a\mathfrak b\mathfrak c=\mathfrak m\\
                   N\mathfrak a,N\mathfrak b\le T}}
                                         \mu(\mathfrak c).
\end{equation}
The factors are pairwise coprime and squarefree. For $Y\le U$ the
difference of cutoff boxes has two disjoint pieces. After interchanging
$\mathfrak a,\mathfrak b$ in the second and setting
$\mathfrak z=\mathfrak b\mathfrak c$, the exact coefficient is
\begin{equation}\label{gtl:T-coefficient}
 T_U(\mathfrak z)+T_Y(\mathfrak z),\qquad
 T_L(\mathfrak z)=
       \sum_{\substack{\mathfrak b\mid\mathfrak z\\N\mathfrak b\le L}}
                                  \mu(\mathfrak z/\mathfrak b).
\end{equation}
It is independent of $\mathfrak a$. The conjugate Gauss CRT identity
retains the cross interaction
\begin{equation}\label{gtl:gauss-cross}
 \overline{G_{\mathfrak a\mathfrak z}}
  =\overline{G_{\mathfrak a}}\overline{G_{\mathfrak z}}
                 \chi_{3,\mathfrak a}(\alpha_{\mathfrak z}),
 \qquad(\mathfrak a,\mathfrak z)=1 .
\end{equation}
The raw cubic entry is zero at a common prime.

The deep input is explicitly external:
\href{https://arxiv.org/abs/2610.04045v1}{de Faveri,
arXiv:2610.04045v1, Theorem~1.1 and Example~2.2}, at $n=3$.
For the physical primary cubic matrix on good squarefree ideals
of norms comparable to $A,Z$ it gives
\begin{equation}\label{gtl:cubic-input}
 \bigl\|[\chi_{3,\mathfrak a}(\alpha_{\mathfrak z})]\bigr\|_{\rm op}^2
  \ll_\epsilon (AZ)^\epsilon
           \{A+Z+(AZ)^{2/3}\}.
\end{equation}
Both indices are actual ideals and hence cube-free, not arbitrary
primitive-conductor labels. This is the same pinned input as the
cubic comparison in Proposition~\ref{cbk:cubic}.

\begin{theorem}[Actual Gauss-weighted cutoff tail]
\label{gtl:tail-theorem}
For any bounded multiplicative ideal function $\lambda$ on good
ideals, with $|\lambda|\le1$, and $1\le Y\le U\le\sqrt D$,
\begin{equation}\label{gtl:linear-tail}
 \begin{split}
 \left|\sum_{\mathfrak m\ {\rm good,SF}}
  \frac{\mu(\mathfrak m)(A_U-A_Y)(\mathfrak m)
          \overline{G_{\mathfrak m}}\lambda(\mathfrak m)}
       {\sqrt{N\mathfrak m}}W(N\mathfrak m/D)\right|\\
 \ll_\epsilon\mathcal M(W)D^\epsilon
                   \{D^{1/3}+\sqrt{D/Y}\}.
 \end{split}
\end{equation}
The constant is uniform in the conductor or raw deletion set in
$\lambda$. For the actual pure leading model, with every original
finite $K$, class/ray/root/$S$ component and coprime pair retained,
\begin{equation}\label{gtl:finite-tail}
 \begin{split}
 |\mathcal L_0[A_U,A_U]-\mathcal L_0[A_Y,A_Y]|
 \ll B_W\mathcal M(W)D^\epsilon A_1X^{1/6}
                       \{Y^{-1/2}+D^{-1/6}\}.
 \end{split}
\end{equation}
The original unequal-norm projection is also allowed.
\end{theorem}
This uses the cubic interaction before the short-divisor triangle.
It does not assert pointwise or coefficient-norm shrinkage of $A_U$,
nor a separate cancellation theorem for $\mu(\mathfrak c)$.

\subsection{What has been reduced, and what remains}
\label{gtl:finite}
Combining Proposition~\ref{gtl:model} and
Theorem~\ref{gtl:tail-theorem} gives
\begin{equation}\label{gtl:combined}
 \begin{split}
 |\mathcal M_{\rm top}-\mathcal L_0[A_Y,A_Y]|
 &\ll (DX)^\epsilon\left\{
 B_W^2\frac{A_1^2X^{7/6}}{D^2}\right.\\
 &\qquad\left.+
 B_W\mathcal M(W)A_1X^{1/6}
                      (Y^{-1/2}+D^{-1/6})\right\}.
 \end{split}
\end{equation}
At $A_1\asymp H=D^{3/4}$, $X=D^{3/8}$ and
$Y=\min\{U,D^{1/8}\}$, the errors are target-sized. If $Y=U$ the
coefficient tail is exactly zero; otherwise the three displayed
powers are $-1/16$, $3/4$ and $31/48$. The remaining exact coefficients
satisfy
\begin{equation}\label{gtl:long-residual}
 \mu(\mathfrak m)A_Y(\mathfrak m)
  =\sum_{\substack{\mathfrak a\mathfrak b\mathfrak c=\mathfrak m\\
                   N\mathfrak a,N\mathfrak b\le Y}}\mu(\mathfrak c),
 \qquad N\mathfrak c\ge D/Y^2\ge D^{3/4}.
\end{equation}
The small-short-divisor, long residual-M\"obius cubic Gauss bilinear
moment $\mathcal L_0[A_Y,A_Y]$ is still unestimated at the target.


\subsection{The exact discrete top prefix, with every ray retained}
\label{gtl:discrete}
There is also a direct reduction that does not introduce the
principal-density split. Its exact support is
\begin{equation}\label{gtl:discrete-support}
 \mathcal T_X=
 \{(\mathfrak w,\mathfrak s):N\mathfrak w<64,\ \mathfrak s\ {\rm good},
       \ V<N\mathfrak s\le2V,\ N\mathfrak w(N\mathfrak s)^6\le X\}.
\end{equation}
There is no squarefree restriction on $\mathfrak s$.
Overlap of $\mathfrak w,\mathfrak s$ and of
$\mathfrak k,\mathfrak s$ remains allowed.
Let $\widetilde\Theta$ be the same outer set as $\theta$ above,
but with \emph{every} original $\rho$, and put
\begin{equation}\label{gtl:discrete-coefficient}
 \begin{aligned}
 \widetilde d_\theta[f_1,f_2]
 &=\frac{A_1}{D}h_{AB\rho}(\mathcal O)
       c_{f_1}(\mathfrak m)\overline{c_{f_2}(\mathfrak n)}
          \frac{\epsilon(\eta)}{\sqrt F}\\
 &\quad\cdot\frac{\mu(\mathfrak e)\eta(\mathfrak e)}{N\mathfrak e}
                                  \overline{\eta(\mathfrak w)} .
 \end{aligned}
\end{equation}
This is an explicit definition, not division by a density or a
primitive value. With the original $\mathfrak k$ support as in
\eqref{gtl:deletions}, define the normalized exact piece
\begin{equation}\label{gtl:direct-prefix}
 \begin{aligned}
 \mathcal P_X[f_1,f_2]
 &=\sum_{\theta\in\widetilde\Theta}\widetilde d_\theta[f_1,f_2]
     \sum_{\mathfrak k}
           \frac{\mu(\mathfrak k)\eta(\mathfrak k)^2}
                {(N\mathfrak k)^2}\\
 &\quad\cdot
 \sum_{\substack{\mathfrak s:\;(\mathfrak w,\mathfrak s)\in\mathcal T_X}}
   \bar\rho(\mathfrak s)^6
       \mathbf1_{(\mathfrak s,\mathfrak m\mathfrak n)=1}
       \dot G\left(\frac{\lambda_\theta(N\mathfrak s)^6}
                              {(N\mathfrak k)^2}\right).
 \end{aligned}
\end{equation}
Then $\mathcal P_X[A_U,A_U]$ is the actual original top
primitive-prefix contribution divided by $D$.
It is a signed representation-fixed piece, not a positive energy
or the whole infinite dual sum.

Let $\mathcal L_X^K[f_1,f_2]$ be exactly
\eqref{gtl:direct-prefix} with $\dot G$ replaced by $g_0$.
Only in that leading model can the finite $\mathfrak k$ sum
be factored as $Q_2^{T_q}(K;\eta^2)$.
The discrete $\mathfrak s$ sum and its endpoints remain.
Unlike $\mathcal L_0$, this model has no
$\kappa_S g(\mathfrak m)g(\mathfrak n)$ density, no
$\rho^6=1$ projection and no integral replacing the
$\mathfrak s$ family.

\begin{proposition}[Direct original-prefix cutoff reduction]
\label{gtl:direct}
Put $\Lambda=A_1X/D^2$. With precisely the original support,
\begin{equation}\label{gtl:direct-model-error}
 |\mathcal P_X[A_U,A_U]-\mathcal L_X^K[A_U,A_U]|
       \ll B_W^2(DX)^\epsilon A_1V\Lambda ,
\end{equation}
and
\begin{equation}\label{gtl:direct-model-tail}
 \begin{split}
 |\mathcal L_X^K[A_U,A_U]-\mathcal L_X^K[A_Y,A_Y]|
 &\ll B_W\mathcal M(W)(DX)^\epsilon A_1V\\
 &\quad\cdot\{Y^{-1/2}+D^{-1/6}\}.
 \end{split}
\end{equation}
In particular the actual original piece satisfies
\begin{equation}\label{gtl:direct-bound}
 \begin{split}
 |\mathcal P_X[A_U,A_U]-\mathcal L_X^K[A_Y,A_Y]|
 &\ll (DX)^\epsilon A_1V
       \{B_W^2\Lambda\\
 &\qquad+B_W\mathcal M(W)(Y^{-1/2}+D^{-1/6})\}.
 \end{split}
\end{equation}
The original unequal-norm projection is allowed.
\end{proposition}
At the critical pure parameters used above the three powers
are again $-1/16$, $3/4$ and $31/48$, with exact zero
coefficient tail when $Y=U$. This is a reduction of the
\emph{actual discrete} piece to a different, still unestimated,
small-short-divisor model. That model retains every $\rho$,
every nonunit $S$ mode, the exact $\mathfrak w,\mathfrak s$
support, root/conductor data, original finite $K$ and both
long residual-M\"obius column coefficients.

\paragraph{Separate open errors and scope.}\label{gtl:boundary}
The accepted fixed-ray disk count gives the distinct top-band
certificate
\begin{equation}\label{gtl:ray-error}
 \frac{|\mathcal P_{\rm top}-M_{\rm top}|}{D}
 \ll B_W^2(DX)^\epsilon A_1X^{1/18}
                       \{1+\sqrt{A_1X/D^2}\}.
\end{equation}
Its pure critical power is $37/48>3/4$; this is an upper certificate,
not a lower bound for the actual error. The direct reduction in Proposition~\ref{gtl:direct} does not
introduce that split: its discrete model retains those effects.
No improved bound for the separate discrepancy, for the remaining
small-cutoff models, or for the other original sectors is proved here. The result is an application of a named external
cubic sieve to an actual arithmetic tail, not a new sieve, a complete
analytic Lean formalization, a determinant/zero-location bridge or
a proof of GRH. Fixed-profile constants and the explicit
$\mathcal M(W)$ cannot be replaced by a uniform supnorm-only claim
for varying profiles.
